\begin{lemma}
\label{lemma-generalized-valuation-ring-modules}
\begin{reference}
\cite[Theorem 1]{Warfield-Decomposition}
\end{reference}
Let $R$ be a ring satisfying the equivalent conditions of
Lemma \ref{lemma-generalized-valuation-ring}.
Then every finitely presented $R$-module
is isomorphic to a finite direct sum of modules of the form $R/fR$.
\end{lemma}

\begin{proof}
Let $M$ be a finitely presented $R$-module. We will use all
the equivalent properties of $R$ from
Lemma \ref{lemma-generalized-valuation-ring}
without further mention. Denote $\mathfrak m \subset R$
the maximal ideal and $\kappa = R/\mathfrak m$ the residue field.
Let $I \subset R$ be the annihilator of $M$.
Choose a basis $y_1, \ldots, y_n$ of the finite dimensional
$\kappa$-vector space $M/\mathfrak m M$. We will argue
by induction on $n$.

\medskip\noindent
By Nakayama's lemma any collection of elements $x_1, \ldots, x_n \in M$
lifting the elements $y_1, \ldots, y_n$ in $M/\mathfrak m M$
generate $M$, see Algebra, Lemma \ref{algebra-lemma-NAK}.
This immediately proves the base case $n = 0$ of the induction.

\medskip\noindent
We claim there exists an index $i$ such that for any choice
of $x_i \in M$ mapping to $y_i$ the annihilator of $x_i$ is $I$.
Namely, if not, then we can choose $x_1, \ldots, x_n$ such
that $I_i = \text{Ann}(x_i) \not = I$ for all $i$. But as
$I \subset I_i$ for all $i$, ideals being totally ordered
implies $I_i$ is strictly bigger than $I$ for $i = 1, \ldots, n$,
and by total ordering once more we would see that
$\text{Ann}(M) = I_1 \cap \ldots \cap I_n$ is bigger than $I$
which is a contradiction. After renumbering we may assume that
$y_1$ has the property: for any $x_1 \in M$ lifting $y_1$
the annihilator of $x_1$ is $I$.

\medskip\noindent
We set $A = Rx_1 \subset M$. Consider the exact sequence
$0 \to A \to M \to M/A \to 0$. Since $A$ is finite, we see that
$M/A$ is a finitely presented $R$-module
(Algebra, Lemma \ref{algebra-lemma-extension}) with fewer generators.
Hence $M/A \cong \bigoplus_{j = 1, \ldots, m} R/f_jR$ by induction.
On the other hand, we claim that $A \to M$ satisfies the property:
if $f \in R$, then $fA = A \cap fM$. The inclusion $fA \subset A \cap fM$
is trivial. Conversely, if $x \in A \cap fM$, then $x = gx_1 = f y$
for some $g \in R$ and $y \in M$. If $f$ divides $g$, then $x \in fA$
as desired. If not, then we can write $f = hg$ for some $h \in \mathfrak m$.
The element $x'_1 = x_1 - hy$ has annihilator $I$ by the previous
paragraph. Thus $g \in I$ and we see that $x = 0$ as desired.
The claim and Lemma \ref{lemma-characterize-PD-modules}
imply the sequence $0 \to A \to M \to M/A \to 0$ is split
and we find $M \cong A \oplus \bigoplus_{j = 1, \ldots, m} R/f_jR$.
Then $A = R/I$ is finitely presented (as a summand of $M$)
and hence $I$ is finitely generated, hence principal.
This finishes the proof.
\end{proof}